% 由 tools/make_appendix_notes.py 自动生成，勿手改
\section*{第 1 章\quad 范畴}
\noindent\textbf{§1.2}\quad 相反范畴的约定极易记反：本书取 $\mathcal C^{\mathrm{op}}(V,W):=\mathcal C(W,V)$，同时把复合次序整体反转，于是 $\mathcal C^{\mathrm{op}}$ 中的态射 $V\to W$ 对应 $\mathcal C$ 中的态射 $W\to V$。逐项核对含下标的复合式时，务必按这一定义验证配平。\par\medskip
\noindent\textbf{§1.2}\quad 原文把光滑函数构造写成 $C^\infty\colon\mathrm{Man}\to\mathrm{Rings}$，但它是反变函子，规范写法应为 $C^\infty\colon\mathrm{Man}^{\mathrm{op}}\to\mathrm{Rings}$；译文已按反变情形校正。\par\medskip
\noindent\textbf{§1.2}\quad 注意 $\mathrm{mat}(n,m):=\mathrm{Mat}_{m\times n}(k)$ 的指标约定：$n$ 是定义域 $k^n$ 的维数、$m$ 是靶 $k^m$ 的维数，故矩阵有 $m$ 行 $n$ 列。若按“$n\times m$ 矩阵”的日常说法记忆，态射的复合方向会立即错位。\par\medskip
\noindent\textbf{§1.5}\quad 原书此处写作 $F_M:=1_XM$，应为 $F_M(X):=1_XM$（把对象 $X$ 映到向量空间 $1_XM$）；译文已按后者处理。注意 $1_XM$ 是 $M$ 的直和项，$\varphi$ 只能作用在与其源、靶匹配的分量上。\par\medskip
\noindent\textbf{§1.5}\quad Yoneda 引理此处只给陈述而未证明；完整形式是对所有 $F\colon\mathcal C^{\mathrm{op}}\to\mathrm{VECT}(k)$ 有自然同构 $\operatorname{Nat}(\mathcal C(-,X),F)\cong F(X)$，证明只需沿恒等态射追踪像。该引理是后文“生成元”与“导出范畴”论证的隐性工具。\par\medskip
\noindent\textbf{§1.6}\quad 这里的全忠实性依赖 $M=MAeA=(Me)A$，即 $e$ 为满幂等元（$AeA=A$）；若 $e$ 不满，$E$ 只能做到全忠实而未必本质满。这是 Morita 理论中最常见的符号疏漏点。\par\medskip
\noindent\textbf{§1.6}\quad Karoubi 完备化的对象是二元组 $(X,e)$，Hom 空间取 $e'\mathcal C(X,X')e$；$X\mapsto(X,1_X)$ 才是全嵌入，不能省略幂等元 $e$ 的左右夹逼，否则复合不封闭。\par\medskip
\noindent\textbf{§1.7}\quad 原书此处提示作 $e_\rho=\tfrac{1}{|G|}\sum_g\operatorname{Tr}\rho(g)\,g$，缺少因子 $\dim\rho$ 且应为 $\rho(g^{-1})$：对非平凡不可约表示它并不满足 $e_\rho^2=e_\rho$。译文按特征标投影的标准公式校正；书中为笔误。\par\medskip
\section*{第 2 章\quad 上同调}
\noindent\textbf{§2.1}\quad 移位同时改了两件事：指标平移 $i\mapsto i+n$，以及微分整体乘以 $(-1)^n$。后者来自把移位看作与集中在次数 $-n$ 的一维空间作张量积、再按 Koszul 规则把奇次运算符移过：每交换一次顺序就出负号，因此恰好是 $(-1)^n$。核对本章后面 $\operatorname{Hom}^\bullet(C[i],D[j])\cong\operatorname{Hom}^\bullet(C,D)[j-i]$ 所说的“适当符号”，以及第三章 $\mu_k(a_1[i_1],\ldots,a_k[i_k])$ 的移位因子时，都要从这一约定出发；漏掉 $(-1)^n$ 是最常见的移位错误。\par\medskip
\noindent\textbf{§2.1}\quad 证明第一步“由于 $C^i$ 是向量空间”是关键：它保证了每个子空间都有直和补，才能把 $C^i$ 拆成 $H\oplus I\oplus R$ 并使 $d\colon R^i\to I^{i+1}$ 同构。换成 $\mathbb Z$-模（例如 $\mathbb Z\xrightarrow{2}\mathbb Z$）这一步就失效，结构定理与后面的拟同构判别引理都不再成立，见练习 2.1。一般的阿贝尔范畴里同样没有这种分裂。\par\medskip
\noindent\textbf{§2.1}\quad 这里出现的 $h^2=0$、$dhd=d$、$hdh=h$ 不是随手选出的恒等式，而是“收缩数据”的典型关系（还配上 $dh$、$hd$ 为投影、$h$ 次数为 $-1$）。第三章的同调扰动引理正是把这组关系抽象为 $\varphi\varphi^\dagger=1_B+h\mu_1^B+\mu_1^Bh$ 等条件，再据此把 $A_\infty$ 结构沿 $h$ 搬运过去；记住这套关系，后面的树形公式才看得懂。\par\medskip
\noindent\textbf{§2.2}\quad 原书此处把末式印成 $(f+dh)\smile g=f\smile g-d(h\smile g)$，符号有误：由 $d(h\smile g)=dh\smile g+(-1)^{|h|}h\smile dg$ 与 $dg=0$ 得 $d(h\smile g)=dh\smile g$，故应为加号（译文已按加号排）。另外注意两处符号的不同来源：Leibniz 律里的 $(-1)^p$ 来自 $f$ 的次数 $p$，而末式的成立只用 $dg=0$，$df=0$ 仅用于保证 $[f]$ 是上同调类。\par\medskip
\noindent\textbf{§2.2}\quad 三个要点：其一，$\langle[u],[v],[w]\rangle$ 一般是一个集合而非单个类（只有当它只含一个元素时才说“良定义”），所以博罗米恩环那个例子的结论写成“含有非零类/只含零类”。其二，两处的 $\pm$ 必须取相容的符号，最终类才闭；不同 $s,t$ 给出的类相差 $\pm$ 恰当项。其三，这一“选 $s,t$ 再取类”的不定性，正是第三章 $A_\infty$ 语言中 $\mu_3$ 只定义到同伦（或到 $\pm$）的原因。\par\medskip
\noindent\textbf{§2.3}\quad “唯一性可用类似论证证明”省掉的正是同伦 $s$ 的构造，其归纳步骤是：在次数 $0$，差 $\bar\psi^0-\bar\varphi^0$ 是闭的且落在 $\ker(Q^0\to N)=\operatorname{im}d_Q$ 中，用 $P^0$ 的自由性提升得到 $s^0$ 使 $\bar\psi^0-\bar\varphi^0=d_Qs^0$；随后考察 $\bar\psi^{i}-\bar\varphi^{i}-s^{i-1}d_P$（$i<0$），它仍是闭的且落在 $\operatorname{im}d_Q$ 中，再提升一次。每一步都依赖 $P^\bullet$ 由自由（或投射）模构成，这就是定理必须对这类分解陈述的原因。\par\medskip
\noindent\textbf{§2.3}\quad $\xi,\eta$ 都是奇次类，而 Ext 环按 Koszul 符号约定分次交换：$\xi\eta=(-1)^{|\xi||\eta|}\eta\xi=-\eta\xi$，所以 $\xi\eta+\eta\xi=0$ 是由次数决定的，与提升的具体取法无关（多项式环情形那条定理中外代数 $\Lambda_n$ 的全部关系都由此而来）。但像“反方向复合给出 $-\rho$ 的提升”“$\xi$ 自乘为零”这类具体结论，仍须把提升按 $K^\bullet$ 的分次逐块相乘、逐块核对符号才能得到，不能只凭分次交换性断言。\par\medskip
\section*{第 3 章\quad 高阶积}
\noindent\textbf{§3.1}\quad 假设 $\mathrm{char}\,k=2$ 时 $+1=-1$，所以下面三条恒等式里的所有符号都消失了；在特征 $0$ 的域上，须改用带 Koszul 符号的 $b_k$（见 §3.1.4），三条恒等式才严格成立。切勿把特征 $2$ 下的“统一解释”直接搬到特征 $0$。\par\medskip
\noindent\textbf{§3.1}\quad 这里 $|\cdot|$ 是元素在 $A$ 中的次数，不是移位后在 $A[1]$ 中的次数，所以每越过一个输入就贡献 $(|a|-1)$；若误把输入当成 $A[1]$ 的元素（次数 $|a|-1$）再算，会多出符号。另一处常见混淆是 $\mu$ 与 $m$ 两套约定：$m_k=\sigma^{-1}\circ b_k\circ\sigma^{\otimes k}$ 不带额外符号，而 $\mu$ 由右作用 $\bigl((a_1[1]\otimes\cdots\otimes a_k[1])b_k\bigr)[-1]$ 定出，两者只差整体符号约定，切勿混用。\par\medskip
\noindent\textbf{§3.2}\quad 书中只说“每个节点表示某个映射 $b_i$”，未限定 $i\geq2$。若允许 $b_1$ 作节点标记，$k$ 个叶子的树可以带任意长的 $b_1$ 单链，$\bar b_k$ 的定义就不再是有限和（形式的 $\partial_h$ 也需要幂级数收敛性才能逐项求导）。标准约定只取 $i\geq2$：每个内节点至少 $2$ 个儿子，故含 $k$ 个叶子的树至多 $k-1$ 个内节点，求和确为有限和。\par\medskip
\noindent\textbf{§3.2}\quad 这两式都略去了下标与求和限制。$\mu_n^A$ 的式子右端用 $b$ 的 $\mu_r^B$ 与输入端的 $\varphi$，此时 $\varphi$ 已由下式先构造好，故 $r=1,\ i_1=n$ 也可；但 $\varphi_n$ 的递推必须取 $r\geq2$（即每个 $\varphi_{i_j}$ 的下标 $i_j\leq n-1$），否则右端会出现 $\mu_1^B\varphi_n$ 自身造成循环。书中把两式写成同一形式，读者需自行补上这一限制。\par\medskip
\noindent\textbf{§3.3}\quad 普通复形只要求 $\mu_1(\delta)=0$（即 $\delta^2=0$），这里却要求整条 MC 方程：$A_\infty$ 范畴的 Hom 空间自带高阶乘积，$\delta$ 作为“内部的微分”必须与它们相容，$\mu_2(\delta,\delta),\mu_3(\delta,\delta,\delta),\ldots$ 正是度量这种相容性的项；当 $\mathcal C^\bullet$ 退化为只有 $\mu_1$ 的 dg-范畴时 MC 就还原为 $\delta^2=0$。下三角条件保证 $\delta_{i_1i_2}\otimes\cdots$ 足够长时为零，故 MC 的无穷和其实是有限和；去掉下三角（曲情形）就必须另加“$\mu_k=0$ 对充分大的 $k$ 成立”或用完备环才能定义 $\widetilde\mu_k$。\par\medskip
\noindent\textbf{§3.3}\quad 原文（见原书 p.56）此式右端末尾作 $f_k$，与左端只有 $f_1,\ldots,f_{k-1}$ 以及态射列只到 $f_{k-1}$ 都不符，译文已改作 $f_{k-1}$。但按因子个数核对仍有问题：左端 $k-1$ 个因子，右端把相邻的 $f_{i-1},f_i,f_{i+1}$ 换成 $\widetilde f_{i-1},\widetilde f_{i+1}$ 两个因子，只剩 $k-2$ 个。正确的理解是几何式的：$\widetilde f_{i+1}$ 把 $A_{i+2}$ 送进 $\operatorname{Cone}(f_i)$，$\widetilde f_{i-1}$ 从锥送出，中间在扭曲乘积里插入一次 $\delta_i$ 就还原出 $f_i$；使用时按这一含义核对，不必逐字依赖原式的指标。\par\medskip
\noindent\textbf{§3.4}\quad 原书用带下划线的 $\underline{\mathcal C}^\bullet$ 记预范畴、不带下划线的 $\mathcal C^\bullet$ 记它的范畴化，译文两处都排成了 $\mathcal C^\bullet$，于是“称 $\mathcal C^\bullet$ 为 $\mathcal C^\bullet$ 的范畴化”看起来是同义反复。阅读时请记住：本条中前一个 $\mathcal C^\bullet$ 指新构造出的那个范畴（对象集与横截结构来自预范畴），后一个才是原来的预范畴。\par\medskip
\section*{第 4 章\quad 箭图}
\noindent\textbf{§4.1}\quad 乘积顺序务必按"从右到左"：$pq$ 表示先走 $q$ 再走 $p$，所以 $a_1\cdots a_k$ 的拼接条件是 $t(a_i)=h(a_{i+1})$（头尾相接）。若按从左到右理解，后面 $P_v(w)=v\mathbb C Qw$、$\mathbb C Q(v,w)=w\mathbb C Qv$ 和 $vw=\delta_{vw}v$ 的下标会整体反过来。\par\medskip
\noindent\textbf{§4.1}\quad $\operatorname{Ext}^1(S_v,S_w)$ 的指标方向最容易写反：它是 $\#\{a:v\leftarrow w\}$，即从 $w$ 指向 $v$ 的箭头数。原因在投射分解 $\bigoplus_{h(a)=v}P_{t(a)}\to P_v$：同调计算给出 $\operatorname{Ext}^1(S_v,S_w)=\operatorname{Hom}(P_{t(a)},S_w)$，只在存在满足 $h(a)=v$、$t(a)=w$ 的箭头时非零。用单箭头箭图 $v_1\to v_2$ 验算：$\operatorname{Ext}^1(S_2,S_1)=\mathbb C$，而 $\operatorname{Ext}^1(S_1,S_2)=0$。\par\medskip
\noindent\textbf{§4.2}\quad 次数公式 $\deg=\bigl(\beta-(\theta_2-\theta_1)\bigr)/\pi$ 中 $\beta$ 与 $\theta_i$ 都以逆时针为正；把某个相位换成等价的 $\theta+2\pi$ 不改变次数，但把线段的定向改取相反代表元会同时给 $\beta$ 与 $\theta$ 加 $\pi$，次数保持不变。实践中最容易出错的是只给其中一个量加 $\pi$，结果整体差一次移位。\par\medskip
\noindent\textbf{§4.2}\quad 何以"偶长弦对象是带对象的极限"：$\lambda\to0$ 时 $J(\lambda,n)\to J(0,n)$ 成为幂零块，$\alpha_01_n+\alpha_1J(0,n)$ 的秩使复形退化为 $2n$ 项的一维链；$\lambda\to\infty$ 时首项与末项互换，得到方向相反的 $S(1,2n)$。因此带对象在参数边界上退化为偶长弦对象；对 $\lambda\in\mathbb C^*$，带对象与弦对象是不同的不可分解对象。\par\medskip
\noindent\textbf{§4.3}\quad 关系 $gXg^{-1}=\zeta X$ 说明 $g$ 并不固定 $X$，而是把 $X$ 的本征值旋转 $\zeta$；因此 $\mathbb C[X]\star G$ 比不变环 $\mathbb C[X]^G=\mathbb C[X^n]$ 大得多，前者的模保留了 $G$ 在纤维上的作用，这正是"轨形"与"商簇"的差别所在。$g^n=1$ 与 $gXg^{-1}=\zeta X$ 相容，因为 $\zeta^n=1$。\par\medskip
\noindent\textbf{§4.3}\quad 粘合数据的关键在于交叠上 $Y=X^{-1}$，所以 $M_X\otimes_{R_X}R_{XY}$ 只是把 $X$ 局部化：$\mathbb C[X]/(X-\lambda)^n$ 在 $\lambda\ne0$ 时仍非零，而在 $\lambda=0$ 时归零。这正是定理 4.7（射影直线上的相干层分类）中 $S(0,n)$ 只能出现在 $X$ 片、$S(\infty,n)$ 只能出现在 $Y$ 片的原因。写 $\varphi_M$ 时勿忘它同时把 $M_X$ 一侧的 $Y$ 按 $Y=X^{-1}$ 重新解释。\par\medskip
\noindent\textbf{§4.4}\quad 旋转 $R$ 在几何侧只是把多边形转过 $2\pi/(n+1)$，但它在表示论侧就是 Auslander--Reiten 平移 $\tau$：这是导出范畴的固有自等价，一般不能由 $\mathrm{mod}\,Q$ 的普通自等价实现。它把 $S_1$ 送到 $P_n[1]$ 而非某个简单模，正说明必须使用导出范畴；这也是图中 $S_1[1]$ 与 $P_n[1]$ 并列出现的原因。\par\medskip
\section*{第 5 章\quad 来自物理的动机}
\noindent\textbf{§5.1}\quad 原文把 $X$ 直接当作可测空间并写出路径积分，这里省略了两点：其一，Feynman 测度 $[d\gamma]$ 一般不存在 $\sigma$-可加版本，需用 Wick 旋转或格点正则化才能严格定义；其二，物理文献通常把期望值归一化为 $\langle F\rangle=Z^{-1}\int F e^{iS/\hbar}[d\gamma]$，原文未除以 $Z$。\par\medskip
\noindent\textbf{§5.2}\quad 这是超对称局域化（localization）的原型：$S=d_QV$ 且 $d_QF=0$ 时，$\hbar$ 依赖被全微分吸收。严格化还需：$d_Q$ 在测度上无反常、$V$ 的边界项可忽略、$Q$ 生成紧群作用；否则 $\int d_Q(FV)e^{iS/\hbar}$ 可能留下边界贡献。\par\medskip
\noindent\textbf{§5.2}\quad 原文的超对称变换印作 $Q_-x=\epsilon\psi_+$，与 $Q_+x=\epsilon\psi_+$ 完全重复，应为 $\psi_-$（译文已改）。另外 Grassmann 积分规则 $\int d\psi=0$、$\int\psi\,d\psi=1$ 即 Berezin 积分，等价于取最高次外形式的系数；这里的 $Z_{\mathrm{SUSY}}$ 是局部化公式的零维原型。\par\medskip
\noindent\textbf{§5.3}\quad 恒等式 $g(X,JY)=\omega(X,Y)$ 的符号依赖 $J$ 与 $\omega$ 的约定：若改用 $g(-,-)=\omega(-,J-)$ 定义度量，则整体符号翻转。Kähler 条件除三者相容外还要求 $d\omega=0$ 且 $J$ 可积，原文此处只列出三种结构。\par\medskip
\noindent\textbf{§5.3}\quad “经典方程蕴含全纯”只对 $Q_A$-闭的场构型成立；真正计数全纯映射时还要取稳定映射紧化并定义虚拟基本类（Gromov--Witten 理论），因为全纯映射可能出现气泡与自同构。原文用“最简单的情形”回避了这些技术点；A 模型不依赖 $J$，是因为 $J$ 只用来写方程，而计数与面积由 $\omega$ 决定。\par\medskip
\noindent\textbf{§5.3}\quad B 模型零次可观测量（全纯函数）构成的环，在物理文献中称为手征环（chiral ring），其乘法由全纯算符的 OPE 给出；原文用 Dolbeault 上同调作为局部可观测量代数，正是这一物理图像的数学替身。\par\medskip
\noindent\textbf{§5.4}\quad 原文此处把矩阵系数印作 $\langle\phi_i,Q(t)\phi_o\rangle$、把复合印作 $Q(U)\circ Q(V)$；按 $Q(U)\colon Q(B_i)\to Q(B_o)$ 的方向，应分别为 $\langle\phi_o,Q(U)\phi_i\rangle$ 与 $Q(V)\circ Q(U)$（译文已改）。此外“作用量可加”还隐含要求沿粘合边界的场迹相容，即拼接处没有跳跃。\par\medskip
\noindent\textbf{§5.5}\quad 这里 $D^\pi$ 表示分裂导出范畴；同调镜像对称断言的是 $A_\infty$ 拟等价，而不只是三角等价。物理上“导出等价给出相同物理”依赖开闭弦对应：开弦范畴决定闭弦理论（Costello 定理）；数学上还需为 $\mathrm{Coh}(X^\vee)$ 选好增强三角结构，才能谈论 $A_\infty$ 范畴。\par\medskip
\section*{第 6 章\quad A 侧}
\noindent\textbf{§6.1}\quad 符号约定需自行统一。原文 p.128 定义 $dh=\iota_{X_h}\omega$，但 p.129 的 $X_h=\sum_i\frac{\partial h}{\partial p_i}\frac{\partial}{\partial q_i}-\sum_i\frac{\partial h}{\partial q_i}\frac{\partial}{\partial p_i}$ 与 $\omega=\sum_i dp_i\wedge dq_i$ 差一负号；按已给的 $X_h$ 反推应为 $dh=\omega(-,X_h)=-\iota_{X_h}\omega$。后文例 6.56（原书 p.156）用 $\dot q_i=-\partial H/\partial p_i$、$\dot p_i=\partial H/\partial q_i$，即后一约定；推导 Floer 条带的 Cauchy--Riemann 方程时必须始终沿用同一约定。\par\medskip
\noindent\textbf{§6.2}\quad “$J$ 一般”只保证模空间在预期维数处光滑，这里省略了两处：一是标记点处的赋值映射也须横截，否则要用含时 Hamilton 扰动或虚拟基本类框架；二是正合性只排除面积为零的泡，还须先证明面积有正下界，才能由 Gromov 紧致性得到有限性。完整细节见 Seidel 书第二部分（文献 218）。\par\medskip
\noindent\textbf{§6.2}\quad $\Delta$ 取的是 Maslov 指数模 $\pi$ 的代表元：公式给出 $\mathbb Z$-次数，模 $2$ 即得先前的 $\mathbb Z_2$-次数。$U$ 只确定到左乘、右乘分别固定 $T_xL_1$、$T_xL_2$ 的正交变换，其行列式为 $\pm1$，故 $\Delta$ 的模 $\pi$ 约化与 $\deg x$ 均与 $U$ 的选取无关；实际计算时务必先固定 $L_1,L_2$ 的定向与分次。\par\medskip
\noindent\textbf{§6.3}\quad 配对问题说明 $CW$ 上的乘积不能像 $CF$ 那样直接计数：$\phi_1^{kh}L_0$ 与 $\phi_1^hL_0$ 只是 Hamilton 同痕而非相等，故先用 $\widetilde\mu_k$ 落到 $CW(\cdot,\cdot;kh)$，再用拟同构 $S_k$ 拉回。原文未写出的是：这些 $S_k$ 还必须满足相干性条件（与 $\widetilde\mu$ 的复合在高阶上相容），否则 $\mu_k=S_k\widetilde\mu_k$ 不满足 $A_\infty$ 关系。\par\medskip
\noindent\textbf{§6.3}\quad 原文 p.164 末作 $\deg X=k-1$，与同页上一行的“$X^i$ 的次数为 $i(k+1)$”矛盾，应为 $k+1$：由 $\Omega=z^kdz=z^{k+1}(dp+idq)$ 可见沿纤维绕行时相位按 $(k+1)$ 倍变化，故 $\deg X=k+1$。\par\medskip
\noindent\textbf{§6.3}\quad 原文 p.165 此处作“$X_0$ 是两个圆的不交并”，疑为笔误：$X_0=\{x\in T^*S^1\mid\|p\|\le1\}$ 是圆柱 $S^1\times[-1,1]$，不交并的两个圆是其边界 $\partial X_0$；紧接着的 $U_0$ 也是 $\partial X_0$ 上的点。\par\medskip
\noindent\textbf{§6.4}\quad 这里已把第一个和式中的指标改为 $\gamma_{j+1}\otimes\cdots\otimes\gamma_{i-1}$：当 $\mu$ 作用于循环区间 $\gamma_i,\ldots,\gamma_k,\gamma_0,\ldots,\gamma_j$ 时，余下恰是 $\gamma_{j+1},\ldots,\gamma_{i-1}$；原文 p.170 印作 $\gamma_{i+j+1}$，疑为笔误。另外 Hochschild 同调在此记作 $HH_\bullet$，与 $CC_\bullet$ 的下标一致。\par\medskip
\noindent\textbf{§6.4}\quad 开闭映射是链映射并不显然：需把“一端为 Reeb 轨道、另一端为若干 Reeb 弦”的圆柱模空间紧化，使其边界项分别对应 Hochschild 微分与辛同调微分；为此要对含时 Hamilton 函数取一般位置，否则会出现伪全纯圆柱的破裂。原文把验证留给文献 6。\par\medskip
\section*{第 7 章\quad B 侧}
\noindent\textbf{§7.1}\quad 原文称 $\mathrm{QCoh}\,X$ 的对象是“所有相干层”，但按记号 $\mathrm{QCoh}$ 取的应是所有拟相干层，相干层只构成其全子范畴 $\mathrm{Coh}\,X$（这正是下一句所说的全子范畴）；原书此处为笔误，勿据此记忆。\par\medskip
\noindent\textbf{§7.1}\quad 光滑性判据的几何含义：$\mathrm{Perf}^{\bullet}X$ 是“由向量丛生成的有限复形”，$D\mathrm{Coh}^{\bullet}X$ 是全体有界复形；$X$ 光滑时二者相等（每个相干层都有有限局部自由分解），$X$ 有奇点时商 $D^b\mathrm{Coh}X/\mathrm{Perf}X$ 非零，它只探测奇点，这正是 §7.2.5 的奇点范畴 $\operatorname{Sing}X$。\par\medskip
\noindent\textbf{§7.2}\quad 这句的对应顺序按本书记号是反的：$MOD^{\bullet}R$ 为全部（右）模、$Mod^{\bullet}R$ 为有限生成模（见第 1 章 §1.5 的约定），故与仿射簇的对应应为 Coh$\leftrightarrow$Mod、QCoh$\leftrightarrow$MOD。原书此处把两个记号对调，对照 §7.1.4 末尾的范畴塔 $\mathrm{Coh}^{\bullet}X\subset\mathrm{QCoh}^{\bullet}X$、$D\mathrm{Mod}^{\bullet}R\subset D\mathrm{MOD}^{\bullet}R$ 即可确认。\par\medskip
\noindent\textbf{§7.2}\quad 微分公式右端的第二个矩阵正是 Koszul 符号：记 $\sigma=\operatorname{diag}(1,-1)$，则右端第二个矩阵等于 $\sigma f\sigma=(-1)^{|f|}f$（$f$ 为奇次时四个分量依次变号），故 $d(f)=d_Qf-(-1)^{|f|}fd_P$。漏掉这个符号 $d^2$ 一般不为零，这是矩阵分解 Hom 复形最易出错之处。\par\medskip
\noindent\textbf{§7.2}\quad 直接验证：由 $\iota_{\partial_r}(dX_s\wedge\omega)=\delta_{rs}\omega-dX_s\wedge\iota_{\partial_r}\omega$ 可知，$d^2\omega$ 中“先插入再外乘”与“先外乘再插入”的两组项逐项相消，只剩 $\sum_rX_rW_r\omega=W\omega$。又 $d$ 的两部分 $\mathbb Z$-次数分别为 $-1$ 与 $+1$，故 $d$ 本身并非 $\mathbb Z$-齐次，只在 $\mathbb Z_2$ 分次下齐次——这正是矩阵分解范畴取 $\mathbb Z_2$-分次的来源。\par\medskip
\noindent\textbf{§7.3}\quad 原书把 Serre 函子的自然同构印作 $\nu_{X,Y}\colon\mathcal C(X,Y)\to\mathcal C(F(Y),X)^\vee$，这与紧接着所说的“$n$ 维光滑射影簇的导出范畴有 Serre 函子 $-\otimes K[n]$”不相容（按该书面公式，$-\otimes K[n]$ 满足的是关于 $F^{-1}$ 的关系）。译文采用标准方向 $\mathcal C(X,Y)\cong\mathcal C(Y,F(X))^\vee$；核对 Serre 对偶时可写成 $\operatorname{Hom}(X,Y)\cong\operatorname{Hom}(Y,X\otimes K[n])^\vee$，方向记反是最常见的错误。\par\medskip
\noindent\textbf{§7.3}\quad 合成顺序不可交换：$\Phi_P=D\pi_{Y*}\circ D(-\otimes P)\circ D\pi_X^*$ 是先沿 $\pi_X$ 拉回到 $X\times Y$、与核 $P$ 张量、再沿 $\pi_Y$ 推前，得到 $\mathrm{Perf}^{\bullet}X\to\mathrm{Perf}^{\bullet}Y$。核所在空间与函子方向的对应是：$P\in D\mathrm{Coh}^{\bullet}(X\times Y)$ 决定 $X\to Y$，把 $P$ 换成它在交换因子后的转置才给出反向函子；$p_2^*\mathcal E\otimes p_1^*(-)$ 的写法只是把拉回与张量换序的等价格式，不能据此调换 $X,Y$。\par\medskip
\noindent\textbf{§7.3}\quad 原书此条（p.212）有两处把 $f_3$ 印成了 $f_2$：一是“$\mathrm{grMF}^{\bullet}(R,f_2)$ 中与 $D\mathrm{Coh}^{\bullet}X$ 等价的子范畴”，二是“不等价于 $\mathrm{grMF}^{\bullet}(R,f_2)$”。按上下文（本条讨论 $f_3=XY(X-Y)$，$f_2$ 的情形已在上一条处理）应为 $f_3$；译文已改。\par\medskip
\section*{第 8 章\quad 镜像对称}
\noindent\textbf{§8.1}\quad 配对 $\langle u,v\rangle=\sum u_iv_i$ 中 $u\in M$ 是“频率”（出现在单项式指数上），$v\in N$ 是“周期”（单参数子群的生成方向）；$\chi_u\circ\lambda_v(t)=t^{\langle u,v\rangle}$ 说两者相乘才给出数。把 $M,N$ 都识别成 $\mathbb Z^n$ 之后极易把指标写反：本节的射线向量住在 $N$，而对偶锥与势函数的单项式指数住在 $M$。\par\medskip
\noindent\textbf{§8.2}\quad “由正规性推出单射”省略的是 Zariski 主定理：正规簇之间的双有理、纤维有限的态射是开浸入；这里轨道分解说明该态射双有理且纤维有限，再加满射即为同构。若 $X$ 不正规，同一个扇给出的映射仍满射，但可能把若干轨道压到同一点，锥与轨道的双射即失效。\par\medskip
\noindent\textbf{§8.3}\quad “阿米巴 $A_t$ 趋近 $H_{\mathrm{trop}}$”的严格含义是：对任意紧集 $K$，$A_t\cap K$ 与 $H_{\mathrm{trop}}\cap K$ 的 Hausdorff 距离当 $|t|\to\infty$ 时趋于零。证明的关键只是初等估计——系数 $t^{k_m}$ 增长最快的单项式在远离折点处严格支配其余项，故交点只能出现在两个指数并列最大的位置。\par\medskip
\noindent\textbf{§8.4}\quad 三个模型名字里的 1D、2D、3D 指复维数：$Y^1_\square$ 是黎曼曲面（实维 2），$Y^2_\square=(T^\vee,W_t)$ 是复曲面上的 Landau--Ginzburg 模型（实维 4），而 $Y^3_\square$ 是 $T^\vee\times\mathbb C^2$ 中由一个方程切出的超曲面，实维为 $4+4-2=6$。写镜像对应时不要混用复维与实维。\par\medskip
\noindent\textbf{§8.4}\quad 环簇情形下“保典”有一个纯组合的判据：细分所用的每条新射线都落在原锥的高度 $1$ 截面上（这里就是 $\square$ 的格点 $(a,b,1)$），于是细分不引入新的典范除子，$K_{X^3_\boxtimes}=\pi^*K_{X^3_\square}$。若射线来自更高层（如 $(a,b,2)$ 方向），解消就会多出未被消去的除子。\par\medskip
\noindent\textbf{§8.4}\quad 原文用 $S_\Delta$ 表示射线向量的集合，而前文（例 8.5 之后）把一维锥的集合记作 $\Delta_1$，二者应是同一个集合，属于记号不统一（见原书 p.231）。又 $\mathcal O(d)$ 的数据是每条射线上的整数 $d_v$，故 $d$ 的自由度就是射线数，在二维时等于 Newton 多边形的边数。\par\medskip
\noindent\textbf{§8.5}\quad 极对偶的几何内容是“交换 Kähler 自由度与复结构自由度”：$\square$ 一侧给出弱 Fano 簇 $V$ 上的椭圆曲线 $E$，$\diamond$ 一侧给出 $V^\vee$ 上的 $E^\vee$，而 $W$ 与 $W^\vee$ 的系数分别充当对方的模参数。二维反射多边形恰有 $16$ 个，对应 $16$ 类 del Pezzo 型弱 Fano 曲面；这是 Batyrev 对偶在复二维的实例。\par\medskip
\noindent\textbf{§8.6}\quad 第二处同构用了 $\Lambda\cong\mathbb Z^n$：逐个生成元提升即可看出 $\operatorname{Hom}(\Lambda,\mathbb R)\to\operatorname{Hom}(\Lambda,\mathbb R/\mathbb Z)$ 满射，其核正是 $\operatorname{Hom}(\Lambda,\mathbb Z)$。最后一步把 $\Lambda$ 换成对偶格 $\Lambda^*=\operatorname{Hom}(\Lambda,\mathbb Z)$；若改变 $\Lambda$ 与 $\pi_1(L_b)$ 的同一化，$B$ 的坐标会差一个 $\mathrm{GL}(n,\mathbb Z)$ 变换，这正是镜像构造中的整变换不确定性。\par\medskip
\section*{第 9 章\quad 粘合}
\noindent\textbf{§9.1}\quad 定义正文说 $\beta_1,\ldots,\beta_l$ 是一个多边形的连续各角，两条乘积公式却写成 $\beta_k$；原书 p.253 此处指标不一致，按上下文两处应指同一组角，代入时以正文的 $\beta_l$ 为准。\par\medskip
\noindent\textbf{§9.1}\quad 第一条的“当且仅当”要两边验证：若存在闭多边形，其各角按规则相乘得到 $\lambda\,1_a$（$\lambda\ne0$），故有非零高阶乘积；反之任何非零高阶乘积所在浸入多边形的边界都由闭多边形拼成。\par\medskip
\noindent\textbf{§9.2}\quad 下标至关重要：$\lambda_i'=G_i^{-1}\lambda_iG_{i-1}$ 中的 $G_{i-1}$ 来自沿曲线搬运矩阵指标时的一次平移；绕一圈相乘后中间项两两相消，剩下的正是共轭 $T^{\lambda'}=G_1^{-1}T^{\lambda}G_1$，与引理第三条“绕行变换共轭”对应。\par\medskip
\noindent\textbf{§9.2}\quad 证明把 $1_c$ 记作 $p$、把 $\lambda_1^{ac}\alpha_{ac}-\lambda_2^{cb}\alpha_{cb}$ 记作 $p^*$，与原书引理 9.24 的命名（$p^*=1_c$，$p=\lambda_1^{ac}\alpha_{ac}-\lambda_2^{cb}\alpha_{cb}$）正好相反；原书 p.264 亦如此。按下一引理的 $p\cdot p^*=\iota^*_{B_1}$ 与“后者次数为 $1$”，应以引理的命名为准，即 $1_c$ 是零次的 $p^*$。\par\medskip
\noindent\textbf{§9.3}\quad 原书 p.274 把 $\mathrm{Gtl}^{\pm}A_2$ 的第二条关系印作 $\beta_1\alpha_1-1_{a_1}$，但 $\alpha_2$ 与 $\beta_1$ 才互为逆：$\alpha_2\beta_1=1_{a_2}$ 必须配 $\beta_1\alpha_2=1_{a_1}$。译文按后者写作 $\beta_1\alpha_2$；若沿用 $\beta_1\alpha_1$，该代数不会 Morita 等价于 $\mathbb C[\alpha_1\beta_1]$。\par\medskip
\noindent\textbf{§9.4}\quad 原书 p.276 这一推导的末行印作 $=-\,\mathbb{1}$，与所要证的 $\mu(\zeta_1,\xi_1,\eta_1)=1_{M_X}$ 矛盾；译文取 $1_{M_X}$。若按原书取 $-1$，则须同步改变 $\mu(\zeta_1,\xi_1,\eta_1)=1_x$ 的符号约定，属原书笔误。\par\medskip
\noindent\textbf{§9.5}\quad 原文把 $\zeta_1,\zeta_2$ 都写作 $M_X\to M_Y$，但紧接着的 $\zeta_{2,\phi}:M_{Y_\phi}\to M_{X_\phi}$ 方向相反；由上一节 $\mu(\zeta_1,\xi_1,\eta_1)=1_x$ 的循环结构，$\zeta_1,\zeta_2$ 中必有一个反向（原书 p.284 笔误），引用时以上标方向公式为准。\par\medskip
\noindent\textbf{§9.6}\quad 直觉上，二聚体对偶就是镜像对称的组合版本：A 侧的弧系与 B 侧 Jacobi 代数的矩阵分解互为镜像，锯齿圈（箭的极大交替路径）在对偶下变成标记点，这正是亏格与标记点数在对偶下互换的来源。\par\medskip
\section*{第 10 章\quad 分次}
\noindent\textbf{§10.1}\quad 平移下标按 $n\pi$ 而不是 $2\pi$：线场是无向的，相位只在模 $\pi$ 意义下有定义，所以 $\vec L[1]$ 是相位加 $\pi$（半圈）而非绕一整圈；$\mathbb Z$ 分次比 $\mathbb Z_2$ 分次多出的正是这一圈内的位置信息。\par\medskip
\noindent\textbf{§10.1}\quad 绕数 $w(\eta,L)=\Delta_L[0,1]/\pi$ 必为整数：线场是 $P T\dot S$ 的截面，其纤维是长为 $\pi$ 的射影直线，沿闭曲线搬运一圈后相位只能相差 $\pi$ 的整数倍。$w\ne0$ 时不存在整体相位函数 $\theta$，这正是该曲线不可分次的原因。\par\medskip
\noindent\textbf{§10.1}\quad 来自向量场的线场其绕数必为偶数：向量场方向以 $2\pi$ 为周期，而线场只以 $\pi$ 为周期，故沿任何闭曲线向量场转过的 $2\pi$ 整数倍对应线场 $w$ 的偶数。配合 $I_m=1-w/2$，这保证向量场给出的指标都是整数；一般线场的指标可以是半整数。\par\medskip
\noindent\textbf{§10.1}\quad 度量可自由选取是关键一步：在多边形内部取平坦度量后，线场相位绕各角的变化只由内角决定，因此 $\sum_i\deg\alpha_i=k-2$ 与外角和恒为 $2\pi$（即内角和 $(k-2)\pi$）是同一件事。\par\medskip
\noindent\textbf{§10.5}\quad 原书 p.307 把两个短正合列都印成 $0\to R/(X)\xrightarrow{Y}R\xrightarrow{X}R/(Y)\to0$；由 $\zeta_2$ 的图，第二个应为 $0\to R/(Y)\xrightarrow{X}R\xrightarrow{Y}R/(X)\to0$。译文已改正，否则无法同时得到 $R/(X)$ 与 $R/(Y)$ 的两种锥，Ext 代数也少两个生成元。\par\medskip
\noindent\textbf{§10.6}\quad 这里的轨道数 $m$ 同时是内部标记点数与锯齿路径数：移位 $u\mapsto u+1+r$ 在 $\mathbb Z_k$ 上的轨道数决定亏格 $g=(k-m)/2$，所以改变 $r$ 就得到不同亏格的镜像曲面；例 10.11 中 $r=3,4$（$k=5$）分别给出 $m=1$ 与 $m=5$，即 $g=2$ 与 $g=0$。\par\medskip
\section*{第 11 章\quad 稳定化}
\noindent\textbf{§11.1}\quad 定义中 $A^{i}(X,Y)$ 指分次 $\mathrm{Hom}$ 复形的第 $i$ 次部分，$\dim$ 是它的普通维数；交错和保证 $\langle X[1],Y\rangle=-\langle X,Y\rangle$，与格罗滕迪克群中 $[X[1]]=-[X]$ 相容。核对符号时可用这条恒等式自检。\par\medskip
\noindent\textbf{§11.2}\quad 条件 $\beta\cdot\theta\le0$ 的方向与 $\theta$ 的选取绑定：把 $\theta$ 换成 $-\theta$ 就把 $\theta$ 稳定换成反稳定。另外模空间非空要求 $\theta\cdot\alpha=0$，所以 $\theta$ 只在 $\alpha$ 的正交补中取值；核对 King 稳定性时先检查这两个约定。\par\medskip
\noindent\textbf{§11.3}\quad 定理中"半稳定 $\mathbb C$-模"应为半稳定 $\mathbb C Q$-模，即 $Q$ 的表示——原书 p.322 漏写了 $Q$，按上下文每个商 $F_i$ 都是 $\mathbb C Q$-模。\par\medskip
\noindent\textbf{§11.3}\quad (SC4) 中的 $X_i$ 在拟同构意义下唯一，这保证 $\phi^{\pm}(X)$ 与心 $\heartsuit_\sigma$ 良定义；切片按 $\pi$ 周期（见 (SC2)），故 (SC4) 里的相位可以整体平移 $\pi$ 的整数倍，实际只用到 $(0,\pi]$ 这一段，心也正是为此取成半开区间。\par\medskip
\noindent\textbf{§11.4}\quad 特殊曲线的相位就是它在平直度量下的方向角：$\int\omega$ 给出的复向量方向恒定，而 $|Z(L)|\le\ell(\vec L)$ 取等号当且仅当曲线处处沿该方向前进。所以这里判断稳定性是纯几何的：把曲线在同伦类中拉直到最短，得到的折线段就给出 Harder--Narasimhan 分解。\par\medskip
\noindent\textbf{§11.4}\quad 平方根 $\sqrt{\varphi}$ 只有在双覆盖 $\widetilde S$ 上才是整体一形式，双覆盖在 $\varphi$ 的零点与极点处分枝；因此中心荷要在覆盖上积分后再取，长度也相应翻倍——例中直线提升后的总中心荷为 $2\pi$、椭圆为 $4\pi$，正是这个原因。\par\medskip
\noindent\textbf{§11.5}\quad 原书 p.335 此处印作 $e^{-i\epsilon}(w_j+ih_j)$，与同一证明前文的 $a_j=h_j+iw_j$ 及图\ref{fig:strip-flat-structure-coordinates} 的坐标标法不一致；若照印，中心荷会整体取共轭。译文按 $e^{-i\epsilon}(h_j+iw_j)$ 统一。\par\medskip
\noindent\textbf{§11.5}\quad 证明的要点是"局部坐标"：每条有限条带两侧标记点之间的复向量 $a_j=h_j+iw_j$ 给出平直结构的局部坐标，而稳定弧的中心荷恰为这两个复点之差，所以 $\iota$ 是局部同胚；再把开、闭与单射性合起来，才得到 $\mathrm{Flat}(S^{\flat})$ 是一个连通分支而非仅仅是开子集。\par\medskip
\section*{第 12 章\quad 变形}
\noindent\textbf{§12.1}\quad 定义中用直积 $\prod_{k\geq0}$ 而非直和，因此 $\mathrm{HC}^{i}(A)$ 允许同时含无穷多个不同输入个数的项；这一约定在 §12.2 的滤过与指数作用中必不可少，否则 $\exp(\phi)$ 展开出的无穷多项会越出复形。\par\medskip
\noindent\textbf{§12.1}\quad 原文此处作 $\exp(a(X)\partial_X)\cdot b(X)=b(e^{a(X)})$；但 $\exp(a\partial_X)$ 是沿向量场 $a\partial_X$ 的时间 $1$ 流，一般给出 $X\mapsto X+a(X)+\cdots$ 的变量代换。只有当 $a(X)=cX$ 为线性时流才恰为 $X\mapsto e^{c}X=e^{a(X)}$，故一般情形下此处疑为原书笔误（见原书 p.341）；后文“轨道对应单项式”应按变量代换而非按 $e^{a}$ 理解。\par\medskip
\noindent\textbf{§12.1}\quad 原文写 $\ell^i\delta\in\mathrm{HH}^{0}(A)$ 作用在 $\mathrm{HH}^{0}(E)$ 上；但按前述分次，$\ell^j$ 位于 $\mathrm{HH}^{-1}$，而括号给出 $[\ell^i\delta,\ell^j]=nj\ell^{i+j}\in\mathrm{HH}^{-1}$，故作用对象疑应为 $\mathrm{HH}^{-1}(E)$（原书 p.344），此处沿原文保留。\par\medskip
\noindent\textbf{§12.3}\quad $C[\ell_m]$ 与 $C[[\ell_f]]$ 的不对称来自 $\mathrm{HH}^{\bullet}(A)=\prod_i\mathrm{HH}^{i}(A)$ 是直积：$\ell_f$ 的各次幂落在不同的经典次数里，故允许作形式幂级数；而 $\ell_m$ 的各次幂落在同一经典次数中，取无穷和会越出直和，只能保留有限线性组合。\par\medskip
\noindent\textbf{§12.4}\quad 此处 $\Lambda_3$ 的下标 $3$ 指生成元 $\xi,\eta,\zeta$ 的个数，而不是分次次数；$p_x^*,p_y^*,p_z^*$ 对应的是 $2$ 次元素 $\eta\zeta,\zeta\xi,\xi\eta$，而 $\iota^*=\xi\eta\zeta$ 对应 $3$ 次元素，它们都住在同一个外代数 $\Lambda_3$ 里，不要误写成 3 次外幂 $\Lambda^3$。\par\medskip
\noindent\textbf{§12.4}\quad 这里 $\hbar$ 的幂次 $1,3,6,10$ 正是三角数，对应三角形面积随边长增长；出现 $\hbar^{10}$ 这样的跳跃并非排版错误。\par\medskip
\noindent\textbf{§12.6}\quad 这里的“对偶二聚体”指 $Q$：题设已令弧系 $A=Q^{\vee}$，故 $A$ 的对偶二聚体是 $Q$，Landau--Ginzburg 模型取 $\operatorname{Jac}Q$；若误取 $\operatorname{Jac}Q^{\vee}$，就会与下文“$Q$ 须锯齿一致”和 $\operatorname{Ext}_{\operatorname{Jac}Q}(S,S)$ 的记号对不上。\par\medskip
